\chapter{Build: A Risk Engine}\label{ch:rc:build-a-risk-engine}

At six in the morning the risk report must be on the desk heads' screens: every trade repriced under
every scenario, the P\&L summed desk by desk, VaR and \omterm{def:rc:market-risk-measures:es}{expected shortfall} at every node, every limit
breach flagged. Between the end-of-day market data and that deadline sits a batch that has to fit its
window every night; in a 2011 survey of banks by McKinsey, average VaR run times ranged from 2 to 15 hours,
and longer under stress. This last chapter builds the firm's \omterm{def:rc:build-a-risk-engine:engine}{risk engine} on top of the pricing library of
Book~5 and the measures of \cref{ch:rc:market-risk-measures}, and asks the question every such engine
faces: how much exactness can the deadline afford?

\section{Architecture: market data, scenarios, revaluation, aggregation}

\begin{definition}[Risk engine, risk data aggregation]\label{def:rc:build-a-risk-engine:engine}
A \emph{risk engine}\index{risk engine} is the system that turns the firm's positions and market data
into its risk measures: it builds scenarios, revalues the positions under them, aggregates the results
along the firm's structure, and reports them against limits. \emph{Risk data aggregation}\index{risk data aggregation}
is, in the Basel Committee's words, defining, gathering and processing risk data according to the bank's
risk reporting requirements to enable it to measure its performance against its risk tolerance or appetite.
\end{definition}

The engine has four stages (\cref{fig:rc:build-a-risk-engine:arch}). A market-data snapshot fixes today's
curves, spots and volatilities. Scenario generators turn history (or stress designs, chapter~22) into
shifts of named \omterm{def:rc:market-risk-measures:factor}{risk factors}. Revaluation turns each position and scenario into a P\&L, through the
pricing library, so that the \omterm{def:rc:build-a-risk-engine:engine}{risk engine} never prices anything itself. Aggregation sums the P\&L along
the \omterm{def:rc:build-a-risk-engine:hierarchy}{risk hierarchy} and computes the measures. The P\&L array, scenarios by trades, is the engine's
central object: every report is a function of it.

\begin{omfigure}
\begin{tikzpicture}[font=\small,box/.style={draw,rounded corners,minimum height=1cm,align=center,text width=2.3cm}]
\node[box,fill=omThm!10] (md) at (0,1.4) {market-data snapshot};
\node[box,fill=omThm!10] (tr) at (0,0) {trades and\\hierarchy};
\node[box,fill=omThm!10] (hi) at (0,-1.4) {history and\\stress designs};
\node[box] (sc) at (3.4,-1.4) {scenario\\generation};
\node[box] (rv) at (3.4,0.7) {revaluation:\\full, grid,\\delta--gamma};
\node[box,fill=omDef!10] (pl) at (6.8,0) {P\&L array\\scenarios $\times$ trades};
\node[box] (ag) at (10.2,0.7) {aggregation\\by node};
\node[box] (rp) at (10.2,-1.0) {VaR, ES, Euler,\\limits, reports};
\node[box,dashed] (pr) at (3.4,2.5) {pricing library\\(Book~5)};
\draw[->,thick] (md) -- (rv);
\draw[->,thick] (tr) -- (rv);
\draw[->,thick] (hi) -- (sc);
\draw[->,thick] (sc) -- (rv);
\draw[<->,thick,dashed] (pr) -- (rv);
\draw[->,thick] (rv) -- (pl);
\draw[->,thick] (pl) -- (ag);
\draw[->,thick] (ag) -- (rp);
\end{tikzpicture}

\omcaption{The \omterm{def:rc:build-a-risk-engine:engine}{risk engine}'s stages. Revaluation calls the pricing library through its fixed interface;
everything downstream works on the array of P\&L by scenario and trade. Schematic.}\label{fig:rc:build-a-risk-engine:arch}
\end{omfigure}

Book~5's library knows market data, instruments, models and engines, and exposes \texttt{price},
\texttt{sensitivities}, \texttt{reprice} and \texttt{price\_batch} over immutable snapshots bumped by
named \omterm{def:rc:market-risk-measures:factor}{risk factors}. This book adds its rates instruments to it: a \omterm{def:rc:curve-construction:pillar}{pillar} zero curve that bumps \omterm{def:rc:curve-construction:pillar}{pillar} by
\omterm{def:rc:curve-construction:pillar}{pillar}, swaps, and swaptions under the normal model of chapter~4, registered with the library's
\texttt{register} so that its generic calls price them like any other instrument.

\omcode{code/firm/riskengine/firm_riskengine.py}{131}{149}{The swaption engine and the registration of the rates instruments with the pricing library.}\label{lst:rc:build-a-risk-engine:register}

\section{Scenario generation}

Each \omterm{def:rc:stress-testing-and-scenarios:historical}{historical scenario} is a set of bumps on named \omterm{def:rc:market-risk-measures:factor}{risk factors}: here the Treasury curve's eight
\omterm{def:rc:curve-construction:pillar}{pillars} (par yields treated as zero rates, a simplification) and the euro--dollar spot rate. The
engine's test market is that of 23 September 2026: rates from 4.49\% at one year to 5.40\% at thirty,
EUR/USD at 1.1411. It builds 250 one-day scenarios from the most recent year and 250 overlapping
ten-day scenarios; the ten-day moves reach 48 basis points on a \omterm{def:rc:curve-construction:pillar}{pillar} and 2.8\% on the exchange rate.
Filtered or stressed scenarios (chapters~21 and~22) plug in at the same place.

\section{Revaluation: full, grid and sensitivity-based}

\begin{definition}[Full revaluation, revaluation grid]\label{def:rc:build-a-risk-engine:reval}
\emph{Full revaluation}\index{full revaluation} reprices every position with its pricing model in every
scenario. A \emph{revaluation grid}\index{revaluation grid} reprices each position on a small grid of shifts
of each \omterm{def:rc:market-risk-measures:factor}{risk factor} it depends on, once, and computes each scenario's P\&L by interpolating on the grid and
adding the factors' contributions.
\end{definition}

\omterm{def:rc:build-a-risk-engine:reval}{Full revaluation} costs one pricing call per trade and scenario; the engine runs it through the
library's \texttt{price\_batch}, which marks a failing trade instead of stopping the run. The grid costs a fixed number per trade
and factor, whatever the number of scenarios; the \omterm{def:rc:market-risk-measures:dg}{delta--gamma approximation} of chapter~21 costs two per
trade and factor. The grid captures each factor's nonlinearity over the whole range of moves; like the
diagonal \omterm{def:rc:market-risk-measures:dg}{delta--gamma approximation}, it drops the cross effects between factors.

\omcode{code/firm/riskengine/firm_riskengine.py}{206}{222}{Grid revaluation: seven points per factor, spanning the largest scenario move.}\label{lst:rc:build-a-risk-engine:grid}

\begin{example}[Cost and error on a 1,000-trade book]\label{ex:rc:build-a-risk-engine:methods}
The test book holds 600 swaps, 200 swaptions and 200 EUR/USD options, 89 of them expiring within a
month (71 sold). \omterm{def:rc:build-a-risk-engine:reval}{Full revaluation} under 250 scenarios takes 251\,000 pricing calls, the grid 50\,200 and
delta--gamma 17\,400. With one-day scenarios, the grid's 99\% VaR is within 1.59\% of \omterm{def:rc:build-a-risk-engine:reval}{full revaluation} at
every node and delta--gamma within 0.74\%. With ten-day scenarios, delta--gamma overstates the short-dated
FX desk's VaR by 18.45\%, and both approximations miss the rate options' VaR by more than 8.6\%
(\cref{fig:rc:build-a-risk-engine:errors}).
\end{example}

\begin{omfigure}
\begin{tikzpicture}
\begin{axis}[width=12cm,height=5.8cm,ybar,bar width=6pt,
  symbolic x coords={Firm,Rates,FX,Swaps,Rate options,FX short,FX long},xtick=data,
  x tick label style={font=\footnotesize,rotate=25,anchor=north east},
  ylabel={10-day VaR error (\%)},
  tick label style={font=\small},label style={font=\small},
  ymin=-5,ymax=20,grid=major,grid style={gray!25},
  legend columns=3,legend style={font=\footnotesize,at={(0.5,1.03)},anchor=south,draw=none,fill=none,/tikz/every even column/.append style={column sep=8pt}},
  table/col sep=comma]
\addplot[fill=omDef!60,draw=omDef] table[x=node,y=grid]{figdata/rates-credit-risk/29-build-a-risk-engine/errors_10d.csv};
\addplot[fill=omThm!60,draw=omThm] table[x=node,y=dg]{figdata/rates-credit-risk/29-build-a-risk-engine/errors_10d.csv};
\addplot[fill=gray!50,draw=gray] table[x=node,y=plan]{figdata/rates-credit-risk/29-build-a-risk-engine/errors_10d.csv};
\legend{grid,delta--gamma,revaluation plan}
\end{axis}
\end{tikzpicture}

\omcaption{Relative error of the ten-day 99\% VaR against \omterm{def:rc:build-a-risk-engine:reval}{full revaluation}, by node of the hierarchy.
Delta--gamma fails on the short-dated options, whose gamma changes over a ten-day move; both approximations
fail on the swaptions, whose value depends on the product of moves at different \omterm{def:rc:curve-construction:pillar}{pillars}; the plan that
revalues the swaptions in full stays within 0.76\% everywhere. Data: the chapter's tutorial, on US Treasury
yields and ECB reference rates.}\label{fig:rc:build-a-risk-engine:errors}
\end{omfigure}

The two failures have different causes (\cref{fig:rc:build-a-risk-engine:fx}). The short-dated FX
options depend on one factor, and their P\&L over a 2.8\% move is far from a parabola: the grid, which
reprices at the move's size, follows it; the Taylor expansion does not. The swaptions depend on the
forward swap rate, a combination of several \omterm{def:rc:curve-construction:pillar}{pillars}, and their convexity lives in the cross products of
\omterm{def:rc:curve-construction:pillar}{pillar} moves that no per-factor method sees.

\begin{omfigure}
\begin{tikzpicture}
\begin{axis}[width=11cm,height=6cm,
  xlabel={10-day EUR/USD move (\%)},ylabel={desk P\&L (USD million)},
  tick label style={font=\small},label style={font=\small},grid=major,grid style={gray!25},
  legend cell align=left,legend style={font=\footnotesize,at={(0.5,0.04)},anchor=south},
  table/col sep=comma]
\addplot[only marks,mark=o,mark size=1.4pt,draw=gray] table[x=move,y=full]{figdata/rates-credit-risk/29-build-a-risk-engine/fx_short.csv};
\addplot[only marks,mark=x,mark size=1.4pt,draw=omDef] table[x=move,y=grid]{figdata/rates-credit-risk/29-build-a-risk-engine/fx_short.csv};
\addplot[only marks,mark=+,mark size=1.4pt,draw=omThm] table[x=move,y=dg]{figdata/rates-credit-risk/29-build-a-risk-engine/fx_short.csv};
\legend{full revaluation,grid,delta--gamma}
\end{axis}
\end{tikzpicture}

\omcaption{P\&L of the short-dated FX desk (mostly sold options) in each ten-day scenario, against the
scenario's EUR/USD move. The grid sits on the \omterm{def:rc:build-a-risk-engine:reval}{full revaluation}; the delta--gamma parabola overstates the
losses of large moves. Data: the chapter's tutorial.}\label{fig:rc:build-a-risk-engine:fx}
\end{omfigure}

\begin{method}[A revaluation plan]\label{met:rc:build-a-risk-engine:plan}
Measure each approximation's error against \omterm{def:rc:build-a-risk-engine:reval}{full revaluation} node by node, on the scenarios that matter
(the longest horizon, the stressed period). Assign each desk the cheapest method within the error budget,
\omterm{def:rc:build-a-risk-engine:reval}{full revaluation} where none is, and re-run the comparison when the book or the scenarios change.
\end{method}

\omcode{code/firm/riskengine/firm_riskengine.py}{245}{256}{A revaluation plan: each trade revalued by the method of its deepest planned node.}\label{lst:rc:build-a-risk-engine:plan}

\section{Aggregation, hierarchies and limits}

\begin{definition}[Risk hierarchy, risk limit]\label{def:rc:build-a-risk-engine:hierarchy}
A \emph{risk hierarchy}\index{risk hierarchy} is the tree along which positions and their risk are
aggregated: trades into books, books into desks, desks into business lines and the firm. A
\emph{risk limit}\index{risk limit} is a ceiling on a risk measure (VaR, a sensitivity, a stress loss) at a
node of the hierarchy, whose breach triggers escalation.
\end{definition}

VaR and ES are not additive: each node's measure comes from its own P\&L vector, the sum of its trades'
columns of the P\&L array. That is why the engine keeps the array, not per-desk numbers. Euler contributions
(chapter~20) split a node's ES among its children so that the parts add up: each child's average loss in the
parent's worst scenarios.

\begin{example}[Ten-day risk by node]\label{ex:rc:build-a-risk-engine:report}
With \omterm{def:rc:build-a-risk-engine:reval}{full revaluation}, the firm's ten-day 99\% VaR is USD~53.53 million and its 97.5\% ES 50.89 million.
The FX business's own ES is 14.63 million, but its Euler contribution to the firm's is only 3.56 million,
against 47.34 million for rates: the firm's tail is a rates tail. The short-dated FX desk's VaR, 14.58 million,
breaches its limit of 12 million, and the engine flags it.
\end{example}

\omcode{code/firm/riskengine/cpp/firm_riskengine.hpp}{49}{72}{The C++20 kernel: VaR, ES and Euler contributions of a node's P\&L vector (the Rust twin is line for line the same).}\label{lst:rc:build-a-risk-engine:kernel}

\begin{dated}{2026-09}{Risk data aggregation}\label{dat:rc:build-a-risk-engine:bcbs239}
The Basel Committee's principles for effective \omterm{def:rc:build-a-risk-engine:engine}{risk data aggregation} and risk reporting (known as
BCBS~239, January 2013) set fourteen principles, eleven for banks (governance, data architecture,
accuracy, completeness, timeliness, adaptability, and five on reports) and three for supervisors. Global
systemically important banks designated in 2011 or 2012 were to comply by January 2016. The Committee's
latest progress report (November 2023) assessed 31 such banks: only two were fully compliant with all the
principles, and no principle was fully implemented across all banks; its newsletter of January 2026 named
data lineage, governance and cross-border data as continuing challenges. The ECB's guide of May 2024
recalled that none of 25 significant institutions reviewed in 2016 fully followed the principles, and set
seven key areas of supervisory expectations.
\end{dated}

\section{Testing a risk engine}

A \omterm{def:rc:build-a-risk-engine:engine}{risk engine} is tested as a chain. The pricing library carries its own acceptance tests. The rates
instruments are tested against closed forms (a swap at the par rate is worth zero, payer minus receiver
swaption is the forward swap). The approximations are tested against \omterm{def:rc:build-a-risk-engine:reval}{full revaluation} on small moves, where
they must agree. The kernel is tested on a fixture shared by its Python, C++20 and Rust versions, and on
the identity that Euler contributions add up to the parent's ES. Counting pricing calls turns the deadline
into a number the tests can check.

\section{Tutorial: one night's run}\label{tut:rc:build-a-risk-engine:tutorial}

\begin{tutorial}
\textbf{Goal.} Run the engine on the 1,000-trade book under one-day and ten-day \omterm{def:rc:stress-testing-and-scenarios:historical}{historical scenarios}; report
VaR, ES and Euler contributions by node; compare full, grid and delta--gamma revaluation for cost and
error; build a revaluation plan. \textbf{End state:} the numbers of
\cref{ex:rc:build-a-risk-engine:methods,ex:rc:build-a-risk-engine:report} and the three charts.
\begin{tutsteps}
\item \textbf{Market and book}: \texttt{market()}, \texttt{book()}.
\item \textbf{Scenarios}: \texttt{scenarios(1)}, \texttt{scenarios(10)}.
\item \textbf{Revaluation}: \texttt{revaluations(h)}, P\&L arrays and pricing calls per method.
\item \textbf{Reports}: \texttt{reports(h)}, \texttt{errors(h)}, \texttt{euler()}, \texttt{limits()};
\texttt{fig\_rc\_riskengine.py}.
\end{tutsteps}
\textbf{What to change next.} Add a parallel-shift grid for the swaptions; use stressed scenarios from
2020 or 2022; run the kernel in C++ on the P\&L array written to disk.
\end{tutorial}

\section{Build: the risk engine}\label{bld:rc:build-a-risk-engine:riskengine}

\begin{build}
\bfield{Purpose} The firm's overnight risk run: scenarios, revaluation through the pricing library,
aggregation along the hierarchy, VaR, ES, Euler contributions and limits.
\bfield{Interface} \texttt{PillarCurve}, \texttt{IRSwap}, \texttt{Swaption}, \texttt{NormalRates}
(registered with the pricing library); \texttt{Trade}; \texttt{historical\_scenarios};
\texttt{full\_revaluation}, \texttt{grid\_revaluation}, \texttt{delta\_gamma\_revaluation},
\texttt{planned\_revaluation}, \texttt{Counter}; \texttt{aggregate}, \texttt{risk\_report},
\texttt{euler\_es}, \texttt{check\_limits}; C++20 and Rust kernels.
\bfield{Rules} Prices come only from the pricing library; the P\&L array is kept; every approximation is
measured against \omterm{def:rc:build-a-risk-engine:reval}{full revaluation}; one failing trade never stops the run.
\bfield{Acceptance tests} \texttt{code/firm/riskengine/tests/}, \texttt{cpp/}, \texttt{rust/}: closed forms of the
rates instruments, \omterm{def:rc:curve-construction:pillar}{pillar} bumps, call counts of each method, agreement on small moves, plans, limits, and the
shared kernel fixture.
\bfield{Stretch} \omterm{def:rc:rates-risk:crossgamma}{Cross-gamma} or parallel-shift grids; incremental runs for new trades; parallel revaluation;
stressed and filtered scenarios; lineage from each number back to its trades and data.
\end{build}

\begin{omsources}
\item Basel Committee on Banking Supervision, \emph{Principles for effective risk data aggregation and risk
reporting}, January 2013.
\item Basel Committee on Banking Supervision, \emph{Progress in adopting the Principles for effective risk data
aggregation and risk reporting}, November 2023.
\item Basel Committee on Banking Supervision, Newsletter No~36 on risk data aggregation, January 2026.
\item European Central Bank, \emph{Guide on effective risk data aggregation and risk reporting}, May 2024.
\item US Department of the Treasury, Daily Treasury Par Yield Curve Rates; European Central Bank, euro foreign
exchange reference rates (data).
\end{omsources}

\section{Exercises}

\begin{exercise}[$\star$]\label{exo:rc:build-a-risk-engine:1}
Count the pricing calls of full, grid (seven points) and delta--gamma revaluation for one swap and one FX
option under 250 scenarios, and check the book's totals.
\end{exercise}

\begin{exercise}[$\star$]\label{exo:rc:build-a-risk-engine:2}
Why does the grid's cost not depend on the number of scenarios?
\end{exercise}

\begin{exercise}[$\star$]\label{exo:rc:build-a-risk-engine:3}
Why must the engine keep the P\&L array rather than each desk's VaR?
\end{exercise}

\begin{exercise}[$\star\star$]\label{exo:rc:build-a-risk-engine:4}
How large is the diversification between rates and FX in the firm's ten-day ES?
\end{exercise}

\begin{exercise}[$\star\star$]\label{exo:rc:build-a-risk-engine:5}
Why does delta--gamma revaluation overstate the short-dated desk's losses on large moves?
\end{exercise}

\begin{exercise}[$\star\star$]\label{exo:rc:build-a-risk-engine:6}
Why does no per-factor method capture the swaptions' ten-day risk, and how would you fix it cheaply?
\end{exercise}

\begin{exercise}[$\star\star\star$]\label{exo:rc:build-a-risk-engine:7}
\emph{Coding.} Count the calls of the revaluation plan and check them against the plan's composition.
\end{exercise}

\begin{exercise}[$\star\star\star$]\label{exo:rc:build-a-risk-engine:8}
\emph{Find the flaw.} ``The grid passed the error test last year, so it stays.''
\end{exercise}

\section{Problem: the six o'clock deadline}

\begin{problem}[{Weekend problem --- the six o'clock deadline}]\label{pb:rc:build-a-risk-engine:1}
The firm's real book holds 50,000 trades like the test book's, the batch may use 16 cores from two to six in
the morning, and one pricing call costs 20 milliseconds on average (assumptions). The risk measure is the
ten-day 99\% VaR under 250 scenarios, and every node's VaR must be within 2\% of \omterm{def:rc:build-a-risk-engine:reval}{full revaluation}.

\medskip\noindent\textbf{Part I --- Cost.}
\begin{enumerate}
\item Give the pricing calls of each method on the test book.
\item Convert them into hours for the real book.
\item Which methods fit the window?
\item What does the grid's cost depend on, and the \omterm{def:rc:build-a-risk-engine:reval}{full revaluation}'s?
\item How many scenarios could \omterm{def:rc:build-a-risk-engine:reval}{full revaluation} afford in the window?
\end{enumerate}

\medskip\noindent\textbf{Part II --- Error.}
\begin{enumerate}[resume]
\item Give the worst node error of each approximation at one day.
\item And at ten days.
\item Which desk breaks delta--gamma, and why?
\item Which desk breaks the grid, and why?
\item Why does the error grow with the horizon?
\end{enumerate}

\medskip\noindent\textbf{Part III --- The plan.}
\begin{enumerate}[resume]
\item Build a plan within the budget from the node errors.
\item Give its cost in calls and in hours.
\item Give its worst node error.
\item Could the plan use delta--gamma for the long-dated FX options? Why is that safe?
\item What must be monitored for the plan to stay valid?
\end{enumerate}

\medskip\noindent\textbf{Part IV --- Judgement.}
\begin{enumerate}[resume]
\item Which limit breaches does the report show?
\item What does the Euler split tell the firm about its tail?
\item What would BCBS~239's timeliness and accuracy principles ask of this run?
\item State the \emph{named result}: the hours and worst VaR error of full, grid, delta--gamma and planned
revaluation, and the cheapest method within the 2\% budget.
\item Why is the cheapest method within budget not a fixed choice?
\end{enumerate}
\end{problem}

\section{Interview questions}

\begin{interviewq}[$\star$ \iqroles{risk}]\label{iq:rc:build-a-risk-engine:1}
Describe the architecture of a \omterm{def:rc:build-a-risk-engine:engine}{risk engine}.
\end{interviewq}

\begin{interviewq}[$\star\star$ \iqroles{developer}]\label{iq:rc:build-a-risk-engine:2}
How would you make an overnight VaR run on 50,000 trades fit a four-hour window?
\end{interviewq}

\begin{interviewq}[$\star\star$ \iqroles{risk, researcher}]\label{iq:rc:build-a-risk-engine:3}
Compare \omterm{def:rc:build-a-risk-engine:reval}{full revaluation}, grids and sensitivity-based revaluation.
\end{interviewq}

\begin{interviewq}[$\star\star$ \iqroles{risk}]\label{iq:rc:build-a-risk-engine:4}
Why are VaR numbers not additive across desks, and how do you allocate them?
\end{interviewq}

\begin{interviewq}[$\star\star\star$ \iqroles{developer, risk}]\label{iq:rc:build-a-risk-engine:5}
How do you test a \omterm{def:rc:build-a-risk-engine:engine}{risk engine}?
\end{interviewq}

\begin{interviewq}[$\star\star\star$ \iqroles{bank, risk}]\label{iq:rc:build-a-risk-engine:6}
What does BCBS~239 require, and why do banks struggle to comply?
\end{interviewq}
